\section{Turing Machine --- Build --- NDHost}
%%% generated by the blueprint pipeline from TCSlib.Complexity.TuringMachine.Build.NDHost
\subsection*{Turing Machine --- Build --- NDHost: Introduction}

Declarations of the Lean module \texttt{TCSlib.Complexity.TuringMachine.Build.NDHost}.

\subsection*{Turing Machine --- Build --- NDHost: Results}

\begin{definition}[Nondeterministic phase host]
\lean{Turing.NDTM.phaseHostND}
\label{Turing.NDTM.phaseHostND}
\leanok
\uses{Turing.Action.mapState, Turing.NDTM}
Let $I$ be a type, let each $i \in I$ have a state type $\mathrm{Ph}(i)$, and let $M_i$
be a nondeterministic Turing machine with $k$ work tapes over alphabet $\Gamma$ and
state type $\mathrm{Ph}(i)$. Given an initial state $q_0$ in the dependent sum of these
state types and a function that may redirect each phase state to another dependent-sum
state, the phase host has state type $\sum_{i\in I}\mathrm{Ph}(i)$ and initial state
$q_0$. At a state $(i,q)$, if the redirect function returns a state, the host makes no
writes or output, leaves all heads stationary, and moves to that state; otherwise it
performs the transition of $M_i$ for the given choice bit, input symbol, and work-head
symbols, embedding any successor state in phase $i$.
\end{definition}

\begin{theorem}[Fixing a Choice Commutes with Phase Hosting]
\lean{Turing.NDTM.phaseHostND_fixChoice}
\label{Turing.NDTM.phaseHostND_fixChoice}
\leanok
\uses{Turing.NDTM, Turing.NDTM.fixChoice, Turing.NDTM.phaseHostND, Turing.phaseHostTM}
\difficulty{1}
Let $I$ be an index type, let $\mathrm{Ph}(i)$ be a state type for each $i\in I$, and
let $M_i$ be a nondeterministic Turing machine with state set $\mathrm{Ph}(i)$ and a
common tape count and symbol alphabet. Choose an initial state in the dependent sum of
these state types and a reduction map that assigns to each phase state either no new
state or a phase state in that sum. For either choice bit $b$, fixing choice $b$ in the
phase-hosted machine gives the same deterministic machine as phase-hosting the machines
obtained by fixing choice $b$ in each $M_i$.
\end{theorem}

\begin{theorem}[Phase host run behavior]
\lean{Turing.NDTM.phaseHostND_run}
\label{Turing.NDTM.phaseHostND_run}
\leanok
\uses{Turing.Action.mapState, Turing.Cfg, Turing.Cfg.mapState, Turing.NDTM, Turing.NDTM.phaseHostND, Turing.NDTM.runWith, Turing.NDTM.runWith_append, Turing.NDTM.runWith_cons, Turing.NDTM.runWith_mapState_of_agreeOn, Turing.NDTM.runWith_nil, Turing.NDTM.stepWith}
\difficulty{5}
Let $M_i$ be the machine assigned to each phase $i$, let $q_0$ be an initial state in
the disjoint union of the phase state sets, and let $\mathrm{red}$ assign to each phase
state either no redirection or a new phase state. Fix a phase $i$, an input word $x$,
configurations $c,c'$ of $M_i$ on $x$, a choice word $w$, and a state $s'$ such that
running $M_i$ from $c$ under $w$ gives $c'$ with state $s'$. Suppose
$\mathrm{red}(i,s')=q'$ for some phase state $q'$, and that at every prefix of $w$ of
length strictly less than $|w|$, whenever the run of $M_i$ from $c$ under that prefix
has state $s$, one has $\mathrm{red}(i,s)=\varnothing$. Then, for every prefix length
$u\le |w|$, the phase-host machine run from the phase-tagged configuration $c$ under the
first $u$ bits of $w$ agrees with the run of $M_i$ under that prefix, with states tagged
by $i$; moreover, for every bit $b$, its run under $w$ followed by $b$ ends at the
configuration $c'$ with its state replaced by $q'$.
\end{theorem}

\begin{definition}[Nondeterministic writer states]
\lean{Turing.NDWriterState}
\label{Turing.NDWriterState}
\leanok
For a type $Q$ of controller states, the writer state type has one state
$\mathrm{run}(q)$ for each $q\in Q$, representing writing under controller state $q$,
together with a state $\mathrm{done}$ representing the live exit anchor.
\end{definition}

\begin{definition}[Bounded nondeterministic bit writer]
\lean{Turing.ndWriterTM}
\label{Turing.ndWriterTM}
\leanok
\uses{Turing.NDTM, Turing.NDWriterState}
Let $k$ be a natural number, let $j$ be one of the $k$ work tapes, let $Q$ be a
controller state type, let $q_0\in Q$, let
$\delta:Q\times\{\mathrm{false},\mathrm{true}\}\to Q$, let
$\mathrm{stop}:Q\to\{\mathrm{false},\mathrm{true}\}$, and let $\mathrm{bounded}$ be a
Boolean flag. The resulting nondeterministic machine over Boolean symbols starts in
$\mathrm{run}(q_0)$. In state $\mathrm{done}$ it stays there with no writes or output
and stationary heads; in state $\mathrm{run}(q)$ it enters $\mathrm{done}$ under the
same idle action if $\mathrm{stop}(q)$ holds or if bounded mode is enabled and tape $j$
is blank at its head. Otherwise, for choice bit $b$, it writes $b$ on tape $j$, moves
that tape's head right, leaves all other heads stationary without writing, emits no
output, and enters $\mathrm{run}(\delta(q,b))$.
\end{definition}

\begin{theorem}[Writing a finite tape block]
\lean{Turing.ndWriterTM_run_block}
\label{Turing.ndWriterTM_run_block}
\leanok
\uses{Turing.Cfg, Turing.FinTM.bufferTape, Turing.NDTM.runWith, Turing.NDWriterState, Turing.ndWriterTM}
Let $k\in\mathbb{N}$, let $j$ be one of the $k$ work tapes, and let $x$ be a binary
input. Suppose a configuration $c$ has live writer state, tape $j$ contains a nonblank
symbol at each cell from its head position through the next $m-1$ cells, and the cell
$m$ positions ahead is blank. For every binary word $w$ of length $m+1$, running the
specified writer machine from $c$ under $w$ ends in its done state; the tape contents
are those of $c$ with the cells from the initial head position up to, but not including,
the cell $m$ positions ahead replaced by the first $m$ bits of $w$, and the head is
advanced by $m$. Moreover, no run from $c$ under a word of length less than $m+1$ is in
the done state.
\end{theorem}

\begin{definition}[Phases of a guess-and-run machine]
\lean{Turing.GuessRunPhase}
\label{Turing.GuessRunPhase}
\leanok
The guess-and-run phase type consists of six phases: deterministic setup, bounded choice
writing, a step to the right past the blank after the guessed block, rewinding the
writer's tape to the block's start, a deterministic run up to a live verdict, and
verdict emission.
\end{definition}

\begin{definition}[State types for each phase]
\lean{Turing.GuessRunPhase.State}
\label{Turing.GuessRunPhase.State}
\leanok
\uses{Turing.GuessRunPhase, Turing.NDWriterState, Turing.SeekPhase}
For state types $S_P$ and $S_Q$, assign to the six phases of a guessing-and-running
computation the following state types: $S_P$ in the setup phase, the two-state
nondeterministic writer type in the guessing phase, the four-state seek type in each of
the skipping and rewind phases, $S_Q$ in the deterministic run phase, and the Boolean
verdict type in the output phase.
\end{definition}

\begin{definition}[Machines for the computation phases]
\lean{Turing.guessRunPhases}
\label{Turing.guessRunPhases}
\leanok
\uses{Turing.GuessRunPhase, Turing.GuessRunPhase.State, Turing.MultiTapeTM, Turing.MultiTapeTM.toNDTM, Turing.NDTM, Turing.ndWriterTM, Turing.skipBackTM, Turing.skipTM}
Let $P$ and $Q$ be deterministic $k$-work-tape Turing machines over the Boolean
alphabet, with state types $S_P$ and $S_Q$, and choose a work tape $j$. The
phase-indexed nondeterministic machines are respectively $P$ with its transition
independent of the choice bit; a bounded writer that writes the chosen Boolean on tape
$j$ until it encounters a blank and then halts; a rightward scan of tape $j$ until a
blank is found; the specified rewind scan of tape $j$; $Q$ with its transition
independent of the choice bit; and a machine that halts after emitting the chosen
Boolean.
\end{definition}

\begin{definition}[Phase transition at designated states]
\lean{Turing.guessRunRed}
\label{Turing.guessRunRed}
\leanok
\uses{Turing.GuessRunPhase, Turing.GuessRunPhase.State}
Let $p_{\rm exit}$ be a setup state, and let $q_{\rm entry}$, $q_{\rm acc}$, and $q_{\rm
rej}$ be states of the deterministic run machine. From the setup phase, the designated
exit state transitions to the initial guessing state, while other states have no phase
transition. The completed writer transitions to the initial skipping state; completion
of the skip transitions to the initial rewind state; and completion of the rewind
transitions to the run phase at $q_{\rm entry}$. In the run phase, $q_{\rm acc}$
transitions to the output phase with verdict true, and otherwise $q_{\rm rej}$
transitions there with verdict false; all other states, including every state in the
output phase, have no phase transition.
\end{definition}

\begin{definition}[Combined guess-and-run machine]
\lean{Turing.guessRunTM}
\label{Turing.guessRunTM}
\leanok
\uses{Turing.GuessRunPhase.State, Turing.MultiTapeTM, Turing.NDTM, Turing.NDTM.phaseHostND, Turing.guessRunPhases, Turing.guessRunRed}
Let $P$ and $Q$ be deterministic $k$-work-tape Turing machines over Booleans, with state
types $S_P$ and $S_Q$. Choose a setup exit state $p_{\rm exit}$, a work tape $j$, and run
states $q_{\rm entry}$, $q_{\rm acc}$, and $q_{\rm rej}$; the nondeterministic machine with
states tagged by computation phase starts in $P$'s initial state in the setup phase and
combines the phase machines with transitions at these designated states: it writes a
bounded guessed block on tape $j$, skips past the following blank, rewinds that tape,
runs $Q$ from $q_{\rm entry}$, and emits true at $q_{\rm acc}$ or false at $q_{\rm rej}$.
\end{definition}

\begin{definition}[Configuration with a guessed tape block]
\lean{Turing.guessRunCfg}
\label{Turing.guessRunCfg}
\leanok
\uses{Turing.Cfg, Turing.FinTM.bufferTape}
Let $c$ be a configuration on Boolean input $z$ with $k$ work tapes and state type $S$,
choose a work tape $j$, a finite Boolean word $s$, and a state $q$ of type $S'$. The
resulting configuration has state $q$, retains the input-head position, all work-head
positions, and the output, and replaces the contents of tape $j$ on the integer interval
from its current head position through the next $|s|$ cells by the word $s$, leaving its
other cells and all other work tapes unchanged.
\end{definition}

\begin{definition}[Initial streaming configuration]
\lean{Turing.streamStartCfg}
\label{Turing.streamStartCfg}
\leanok
\uses{Turing.Cfg, Turing.FinTM.bufferTape}
Let $z$ be a Boolean input, choose one of $K$ work tapes, and let $q$ be a state and $s$
a finite Boolean word. The initial configuration has state $q$, input-head position $1$,
all work-head positions $0$, empty output, and tape contents equal to $s$ on the chosen
tape and blank everywhere on every other tape.
\end{definition}

\begin{theorem}[Composition of state renamings]
\lean{Turing.guessRun_mapState_mapState}
\label{Turing.guessRun_mapState_mapState}
\leanok
\uses{Turing.Cfg, Turing.Cfg.mapState}
\difficulty{3}
Let $c$ be a configuration of a machine with $k$ work tapes on a Boolean input, and let
$f:S\to S'$ and $g:S'\to S''$ be functions between state sets. Renaming the state of $c$
first by $f$ and then by $g$ gives the same configuration as renaming it once by the
composite function $g\circ f$.
\end{theorem}

\begin{theorem}[Splitting a Boolean word at a given position]
\lean{Turing.guessRun_split}
\label{Turing.guessRun_split}
\leanok
\difficulty{3}
Let $w$ be a finite Boolean word and let $L$ be a natural number with $L<|w|$. Then
there are Boolean words $a,r$ and a Boolean $b$ such that $w=a\,b\,r$ and $|a|=L$.
\end{theorem}

\begin{theorem}[Handoff to the verifier phase]
\lean{Turing.guessRunTM_runWith_handoff}
\label{Turing.guessRunTM_runWith_handoff}
\leanok
\uses{Turing.Cfg, Turing.Cfg.mapState, Turing.FinTM.bufferTape, Turing.FinTM.bufferTape_left, Turing.FinTM.bufferTape_nat, Turing.GuessRunPhase.State, Turing.MultiTapeTM, Turing.MultiTapeTM.initCfg, Turing.MultiTapeTM.runFrom, Turing.MultiTapeTM.toNDTM, Turing.MultiTapeTM.toNDTM_runWith, Turing.NDTM.initCfg, Turing.NDTM.phaseHostND, Turing.NDTM.phaseHostND_run, Turing.NDTM.runWith, Turing.NDTM.runWith_append, Turing.NDWriterState, Turing.SeekPhase, Turing.guessRunCfg, Turing.guessRunPhases, Turing.guessRunRed, Turing.guessRunTM, Turing.guessRun_mapState_mapState, Turing.guessRun_split, Turing.ndWriterTM_run_block, Turing.skipBackTM, Turing.skipBackTM_run_ofCfg, Turing.skipTM, Turing.skipTM_run_ofCfg}
\difficulty{7}
Let $P$ and $Q$ be deterministic machines with $k$ Boolean work tapes, let $j$ be a work
tape index, and choose states $p_{\mathrm{exit}}$ of $P$ and
$q_{\mathrm{entry}},q_{\mathrm{acc}},q_{\mathrm{rej}}$ of $Q$. Suppose $P$ reaches
configuration $c_P$ in state $p_{\mathrm{exit}}$ for the first time after $T_P$ steps on
Boolean input $z$. Assume tape $j$ in $c_P$ has a block of $m$ nonblank cells starting
at its head position, followed by a blank, and the cell immediately to the left of the
head is blank. For any Boolean word $w$ of length $T_P+2m+9$, after the guess-and-run
machine processes $w$ from its initial configuration, its configuration is the one
obtained from $c_P$ by replacing that block on tape $j$ with the length-$m$ word formed
by the $m$ bits of $w$ immediately following its first $T_P$ bits, preserving all other
tape contents and positions, and entering state $q_{\mathrm{entry}}$ in the run phase.
\end{theorem}

\begin{theorem}[Liveness before the time bound]
\lean{Turing.guessRunTM_runWith_live}
\label{Turing.guessRunTM_runWith_live}
\leanok
\uses{Turing.Cfg, Turing.Cfg.mapState, Turing.MultiTapeTM, Turing.MultiTapeTM.initCfg, Turing.MultiTapeTM.runFrom, Turing.NDTM.initCfg, Turing.NDTM.runWith, Turing.NDTM.runWith_append, Turing.NDTM.runWith_of_halt, Turing.guessRunCfg, Turing.guessRunTM, Turing.guessRunTM_runWith_handoff}
\difficulty{3}
Let $P$ and $Q$ be deterministic machines with $k$ Boolean work tapes, let $j$ be a work
tape index, and choose states $p_{\mathrm{exit}}$ of $P$ and
$q_{\mathrm{entry}},q_{\mathrm{acc}},q_{\mathrm{rej}}$ of $Q$. Suppose $P$ reaches
configuration $c_P$ in state $p_{\mathrm{exit}}$ for the first time after $T_P$ steps on
Boolean input $z$. Assume tape $j$ in $c_P$ has a block of $m$ nonblank cells starting
at its head position, followed by a blank, and the cell immediately to the left of the
head is blank. For every Boolean word $w$ of length strictly less than $T_P+2m+9$, the
guess-and-run machine's configuration after processing $w$ from its initial
configuration is not halted.
\end{theorem}

\begin{theorem}[State at the verifier handoff]
\lean{Turing.guessRunTM_runWith_handoff_state}
\label{Turing.guessRunTM_runWith_handoff_state}
\leanok
\uses{Turing.Cfg, Turing.MultiTapeTM, Turing.MultiTapeTM.initCfg, Turing.MultiTapeTM.runFrom, Turing.NDTM.initCfg, Turing.NDTM.runWith, Turing.guessRunTM, Turing.guessRunTM_runWith_handoff}
\difficulty{3}
Let $P$ and $Q$ be deterministic machines with $k$ Boolean work tapes, let $j$ be a work
tape index, and choose states $p_{\mathrm{exit}}$ of $P$ and
$q_{\mathrm{entry}},q_{\mathrm{acc}},q_{\mathrm{rej}}$ of $Q$. Suppose $P$ reaches
configuration $c_P$ in state $p_{\mathrm{exit}}$ for the first time after $T_P$ steps on
Boolean input $z$. Assume tape $j$ in $c_P$ has a block of $m$ nonblank cells starting
at its head position, followed by a blank, and the cell immediately to the left of the
head is blank. Then, for every Boolean word $w$ of length $T_P+2m+9$, the configuration
obtained by running the guess-and-run machine on $w$ from its initial configuration has
state equal to the run-phase state with underlying state $q_{\mathrm{entry}}$.
\end{theorem}

\begin{theorem}[Run-and-verdict simulation]
\lean{Turing.guessRunTM_runWith}
\label{Turing.guessRunTM_runWith}
\leanok
\uses{Turing.Action.mapState, Turing.Cfg, Turing.Cfg.mapState, Turing.GuessRunPhase.State, Turing.MultiTapeTM, Turing.MultiTapeTM.initCfg, Turing.MultiTapeTM.runFrom, Turing.MultiTapeTM.toNDTM, Turing.MultiTapeTM.toNDTM_runWith, Turing.NDTM.initCfg, Turing.NDTM.phaseHostND, Turing.NDTM.phaseHostND_run, Turing.NDTM.runWith, Turing.NDTM.runWith_append, Turing.NDTM.runWith_cons, Turing.NDTM.runWith_of_halt, Turing.NDTM.stepWith, Turing.guessRunCfg, Turing.guessRunPhases, Turing.guessRunRed, Turing.guessRunTM, Turing.guessRunTM_runWith_handoff, Turing.guessRun_split}
\difficulty{6}
Let $P$ and $Q$ be deterministic machines with $k$ Boolean work tapes, let $j$ be a work
tape index, and choose states $p_{\mathrm{exit}}$ of $P$ and
$q_{\mathrm{entry}},q_{\mathrm{acc}},q_{\mathrm{rej}}$ of $Q$. Suppose $P$ reaches
configuration $c_P$ in state $p_{\mathrm{exit}}$ for the first time after $T_P$ steps on
input $z$, and tape $j$ at that configuration has a block of $m$ nonblank cells from its
head position, followed by a blank, with a blank immediately to the left of the head.
Let $s$ be a Boolean word of length $m$, and suppose the run of $Q$ from the
configuration formed by placing $s$ on tape $j$ at the head position of $c_P$ and
entering $q_{\mathrm{entry}}$ reaches $q_{\mathrm{acc}}$ or $q_{\mathrm{rej}}$ for the
first time at step $T_Q$. If a choice-bit word $w$ satisfies $(w\text{ with its first
}T_P+1\text{ bits removed, then truncated to }m\text{ bits})=s$ and $T_P+2m+T_Q+11\le
|w|$, then running the guess-and-run machine on $w$ from its initial configuration
yields a halted configuration whose input-head position and work tapes and tape-head
positions are those of the stated $Q$ configuration, and whose output is that
configuration's output followed by the bit indicating whether its state is
$q_{\mathrm{acc}}$.
\end{theorem}

\begin{theorem}[Halting during the setup phase]
\lean{Turing.guessRunTM_runWith_setup_halt}
\label{Turing.guessRunTM_runWith_setup_halt}
\leanok
\uses{Turing.Cfg.mapState, Turing.GuessRunPhase.State, Turing.MultiTapeTM, Turing.MultiTapeTM.initCfg, Turing.MultiTapeTM.runFrom, Turing.NDTM.initCfg, Turing.NDTM.phaseHostND, Turing.NDTM.runWith, Turing.NDTM.runWith_append, Turing.NDTM.runWith_mapState_of_agreeOn_det, Turing.NDTM.runWith_of_halt, Turing.guessRunPhases, Turing.guessRunRed, Turing.guessRunTM}
\difficulty{5}
Let $P$ be a deterministic machine with $k$ Boolean work tapes, and let $Q$ be another
such machine. Choose a state $p_{\mathrm{exit}}$ of $P$, a work tape index $j$, and
states $q_{\mathrm{entry}},q_{\mathrm{acc}},q_{\mathrm{rej}}$ of $Q$. If $P$ is halted
after $T_P$ steps on Boolean input $z$, and has not visited $p_{\mathrm{exit}}$ at any
earlier time, then for every Boolean word $w$ of length at least $T_P$, running the
guess-and-run machine on $w$ from its initial configuration gives exactly the
configuration reached by $P$ after $T_P$ steps, with its state placed in the pre-phase
of the combined machine.
\end{theorem}

\begin{theorem}[Halting bound for the guess-and-run machine]
\lean{Turing.guessRunTM_haltsWithin}
\label{Turing.guessRunTM_haltsWithin}
\leanok
\uses{Turing.Cfg, Turing.MultiTapeTM, Turing.MultiTapeTM.initCfg, Turing.MultiTapeTM.runFrom, Turing.NDTM.HaltsWithin, Turing.guessRunCfg, Turing.guessRunTM, Turing.guessRunTM_runWith}
\difficulty{3}
Let $P$ and $Q$ be deterministic machines with $k$ Boolean work tapes, let $j$ be one of
those tapes, and choose states $p_{\mathrm{exit}}$ of $P$ and
$q_{\mathrm{entry}},q_{\mathrm{acc}},q_{\mathrm{rej}}$ of $Q$. Suppose that on input
$z$, machine $P$ reaches configuration $c_P$ after $T_P$ steps, with state
$p_{\mathrm{exit}}$, and does not visit that state at any earlier time. Suppose tape $j$
in $c_P$ has a block of at least $m$ nonblank cells beginning at its head position,
followed by a blank cell, and that the cell immediately to the left of the head is
blank. If, for every Boolean word $s$ of length $m$, machine $Q$ started from the
configuration obtained by placing $s$ on tape $j$ at the head position of $c_P$ and
entering state $q_{\mathrm{entry}}$ reaches either $q_{\mathrm{acc}}$ or
$q_{\mathrm{rej}}$ within at most $T$ steps, with neither verdict visited earlier, then
the nondeterministic machine that runs $P$, guesses and writes a bounded block, rewinds,
and then runs $Q$ halts on every choice-bit word of length $T_P+2m+T+11$.
\end{theorem}

\begin{theorem}[Acceptance criterion for the guess-and-run machine]
\lean{Turing.guessRunTM_accepts_iff}
\label{Turing.guessRunTM_accepts_iff}
\leanok
\uses{Turing.Cfg, Turing.MultiTapeTM, Turing.MultiTapeTM.initCfg, Turing.MultiTapeTM.runFrom, Turing.NDTM.initCfg, Turing.NDTM.runWith, Turing.guessRunCfg, Turing.guessRunTM, Turing.guessRunTM_runWith}
\difficulty{6}
Let $P$ and $Q$ be deterministic machines with Boolean work tapes, state sets $SP$ and
$SQ$, and let $pexit$ be a state of $P$, $j$ a work-tape index, and $qentry$, $acc$, and
$rej$ states of $Q$. Fix a finite Boolean input $z$, a time $TP$, and a configuration
$cP$ such that running $P$ from its initial configuration on $z$ for $TP$ steps yields
$cP$, whose state is $pexit$, and no earlier configuration has state $pexit$. Suppose
that, for some natural number $m$, the $m$ cells starting at the head position of tape
$j$ in $cP$ are all nonblank, the next cell is blank, and the cell immediately to their
left is blank. Let $T$ be a natural number such that for every Boolean word $s$ of
length $m$, machine $Q$, started in state $qentry$ with $s$ written on tape $j$ as
prescribed by the guessed-run construction, reaches either $acc$ or $rej$ within some
time $TQ\le T$, and reaches neither verdict at any earlier time. For every natural
number $B$ with $TP+2m+T+11\le B$, there is a choice word of length $B$ that makes the
guess-and-run machine halt with output $[\mathrm{true}]$ if and only if there is a
Boolean word $s$ of length $m$ and a time $TQ$ at which $Q$ reaches $acc$, reaches
neither $acc$ nor $rej$ earlier, and has empty output at time $TQ$.
\end{theorem}

\begin{theorem}[Clean configuration for a guessed run]
\lean{Turing.guessRunCfg_clean}
\label{Turing.guessRunCfg_clean}
\leanok
\uses{Turing.Cfg, Turing.FinTM.bufferTape, Turing.MultiTapeTM, Turing.MultiTapeTM.runFrom_zero, Turing.embedSilentCfg, Turing.embedSilentCfg_cap, Turing.embedSilentCfg_selected_pos, Turing.embedSilentCfg_selected_tape, Turing.embedSilentTM_frame, Turing.guessRunCfg, Turing.streamStartCfg}
\difficulty{6}
Let $P$ and $Q$ be deterministic machines with state sets $SP$ and $SR$, respectively,
over Boolean work tapes, and let $z$ be a finite Boolean input. Suppose $\iota$ embeds
the $K$ tape indices into the $k$ tape indices, the tape index $cap$ lies outside the
image of $\iota$, and $jR<K$. Let $blk$ be a tape content that is blank at every integer
position $p<0$ or $p\ge m$, and let $s$ be a Boolean word of length $m$. For every state
$q$ of $P$ and entry state $entry$ of $Q$, applying the guessed-run configuration
construction to the configuration with state $q$, input-head position $1$, all tape
contents blank except for content $blk$ on tape $\iota(jR)$, all work-tape heads at
zero, and empty output, yields exactly the embedded silent configuration of the
stream-start configuration for $Q$ with state $entry$, tape $jR$ containing $s$, and all
other tapes blank.
\end{theorem}

\begin{theorem}[Silent simulation and acceptance specification]
\lean{Turing.guessRunTM_silent_spec}
\label{Turing.guessRunTM_silent_spec}
\leanok
\uses{Turing.Cfg, Turing.MultiTapeTM, Turing.MultiTapeTM.initCfg, Turing.MultiTapeTM.runFrom, Turing.NDTM.HaltsWithin, Turing.NDTM.initCfg, Turing.NDTM.runWith, Turing.embedSilentCfg, Turing.embedSilentTM, Turing.embedSilentTM_runFrom, Turing.guessRunCfg, Turing.guessRunCfg_clean, Turing.guessRunTM, Turing.guessRunTM_accepts_iff, Turing.guessRunTM_haltsWithin, Turing.streamStartCfg}
\difficulty{5}
Let $K,k$ be natural numbers, let $\iota$ be an injection from the $K$ work-tape indices
into the $k$ work-tape indices, and choose a tape index $\mathrm{cap}$ outside the image
of $\iota$. Let $j_R$ be a work-tape index for a deterministic machine $R$ with $K$
Boolean work tapes, and choose states $\mathrm{entry},\mathrm{acc},\mathrm{rej}$ of $R$.
Let $P$ be a deterministic machine with $k$ Boolean work tapes and choose a state
$p_{\mathrm{exit}}$ of $P$. Suppose a tape profile $\mathrm{blk}:\mathbb
Z\to\mathrm{Option}(\{0,1\})$ is nonblank exactly at integer positions $p$ satisfying
$0\le p<m$. Assume that after $T_P$ steps on input $z$, $P$ is in state
$p_{\mathrm{exit}}$, its tape contents are blank except for this profile on tape
$\iota(j_R)$, every work-tape head is at the origin, its input head is at position $1$,
and its output is empty; also assume it did not visit $p_{\mathrm{exit}}$ earlier.
Suppose that for every Boolean word $s$ of length $m$, the run of $R$ from its initial
configuration modified to enter $\mathrm{entry}$ and place $s$ on tape $j_R$ reaches
either $\mathrm{acc}$ or $\mathrm{rej}$ within $T$ steps, with neither verdict visited
earlier. Then the combined nondeterministic machine halts for every choice-bit word of
length $T_P+2m+T+11$. Moreover, for every $B\ge T_P+2m+T+11$, there exists a length-$B$
choice-bit word that makes this machine halt with output exactly the one-bit word
$[\mathrm{true}]$ if and only if there is a length-$m$ Boolean word $s$ for which $R$
reaches $\mathrm{acc}$ at some time, having visited neither $\mathrm{acc}$ nor
$\mathrm{rej}$ at any earlier time.
\end{theorem}
